% !PRINTED 667
\weberpage{667}
The definition of the order number is therefore attached to the point $\mathfrak{P}$. It is invariant for the field $\Omega$.

That every function receives at every point $\mathfrak{P}$ a definite \gesperrt{integer} $n$ as its order number now follows from § 185, (5).

Take a function $z$ which remains finite at $\mathfrak{P}$ as independent variable, determine the prime functional $\pi$ belonging to $\mathfrak{P}$, and choose a function $\varrho$ whose numerator is divisible by $\pi$, but not by $\pi^2$. This function is of the first order.

Then every function $\omega$ can be put into the form:

\weq{1}{\omega = a\varrho^n + \eta_1\varrho^{n+1},}
in which $a$ is a constant different from zero, and $n$ is the \gesperrt{order number} of $\omega$.

If $\varrho$ has the order 1, then $\varrho^n$ has the order $n$.

The following theorems hold for the order numbers, and all of them follow easily from (1):

4. \gesperrt{The order of a product of two functions is equal to the sum of the orders of the factors. The order of a quotient is equal to the difference of the orders of numerator and denominator.}

\gesperrt{The order of a sum is equal to the lowest among the orders of the summands. If several summands of the same lowest order number occur, then the order of the sum can be greater, but not smaller.}

\gesperrt{Of a function which becomes zero at $\mathfrak{P}$ in the $n$th order one also says that it becomes zero in the first order at $n$ coincident points.}

\begin{center}
\textbf{§ 189. \gesperrt{Polygons.}}
\end{center}

1. \gesperrt{Definition. Complexes of points, which may also contain the same point several times, are called polygons or many-sided figures} (of the field $\Omega$).

The number of points of a polygon is called its \gesperrt{order}, and a polygon of the $m$th order is also called briefly an $m$-gon.

As designation for the polygons the letters of the capital German alphabet $\mathfrak{A}, \mathfrak{B}, \mathfrak{C}, \ldots$ shall serve.
